%%%%%%%%%%%%%%%%%%%%%%%%%%%%%%%%%%%%%%%%%%%%%%%%%%%%%%%%%%%
% Zusammenfassung des Logbuchs - Björn Becker
% KI 1 Projekt, WS 2025/26 - Gruppe 308
%%%%%%%%%%%%%%%%%%%%%%%%%%%%%%%%%%%%%%%%%%%%%%%%%%%%%%%%%%%

\documentclass[9pt,a4paper,twocolumn,twoside]{tau-class/tau}
\usepackage[ngerman]{babel}
\usepackage{booktabs}
\usepackage{siunitx}
\usepackage{amsmath}
\usepackage{hyperref}

%----------------------------------------------------------
% Title & Author
%----------------------------------------------------------

\doctype{Zusammenfassung des Logbuchs}
\title{Neuronales Netz f\"ur den kalifornischen Hauspreis-Datensatz}

\author[a]{Björn Becker}
\affil[a]{Justus-Liebig-Universität Gießen, Martrikelnummer: \dots}

%----------------------------------------------------------
% Document Info
%----------------------------------------------------------

\dates{Gruppe 308 -- KI~I Projekt WS~2025/26}
\footinfo{KI~I Portfolio}
\theday{25.03.2026}
\organization{JLU Gießen}
\leadauthor{Björn Becker}

%----------------------------------------------------------
\begin{abstract}
    Dieses Dokument fasst meinen Arbeitsfortschritt bei der Entwicklung Neuronaler Netze für die Vorhersage von Hauspreisen anhand des 
    "California housing dataset" zusammen. Es werden vor allem die Hyperparameteroptimierung und Ensemble Erstellung behandelt. 
    Das beste Modell erreicht einen $R^2$-Wert von \dots auf dem Testset und übertrifft damit die lineare Regression ($R^2 = 0{,}63$) deutlich.
\end{abstract}

%----------------------------------------------------------
\begin{document}

\maketitle
\thispagestyle{firststyle}
%----------------------------------------------------------

\section{Fragestellung und Ziel}
Ziel des Projekts war es, ein Neuronales Netz aufzubauen, welches die Hauspreise im "California housing dataset" vorhersagen kann. 
Neuronale Netze stellen eine zentrale Methode der modernen Künstlichen Intelligenz dar (Russell \& Norvig, 2021).
Das Ergebnis sollte mit einer linearen Regression verglichen werden, um den Mehrwert komplexerer Modelle zu bewerten.
Aufgrund der zeitlichen Verfügbarkeit der Gruppenmitglieder hatte ich die Chance die Gruppenergebnisse zusammenzuführen und weiter
zu verwenden. Daher lag der Fokus meiner Arbeit auf der Hyperparameteroptimierung durch die Random-Search-Methode und der Erstellung von Ensembles.

%----------------------------------------------------------

\section{Darstellung der Arbeitsprozesse und Ergebnisse}
\subsection{Vorbereitung}
Vor unserem ersten Gruppentreffen habe ich mich mit der Aufgabe vertraut gemacht und bin die Vorlesungsfolien durchgegangen.
Zusätzlich informierte ich mich über die Einrichtung einer gemeinsamen Programmierumgebung über GitLab. Allerdings kam im Verlauf des 
ersten Gruppentreffens auf, dass insgesamt in der Gruppe mehr Erfahrungen mit GitHub vorhanden waren, weshalb wir uns dafür 
entschieden. Florian übernahm die Einrichtung des Repositories und die Architektur der Ordnerstruktur (vgl. Logbuch von Florian).
\\
Zwischen dem ersten und 5. Gruppentreffen lag mehr als ein Monat, in dem ich lediglich an den Treffen teilnahm, Florian, Felix H. und Kolja bereits den wesentlichen Teil 
ihres Arbeitspensums erledigten. 
Nachdem ich mich einen Tag vor dem fünften Treffen mit dem Fortschritt der Gruppe und den im Repository zugänglichen Dateien, u.a. den Hilfsfunktionen für die Modellevaluation vertraut machte,
machte ich mir Anhand der Vorlesungsfolien Gedanken über mögliche Aufgaben, an denen ich arbeiten könnte. Außerdem habe ich via venv eine 
passende Python-Umgebung eingerichtet und die benötigten Pakete installiert, um die Jupyter-Notebooks der Gruppe ausführen zu können.

\subsection{Erste Modelle}
Nach dem Gruppentreffen, habe ich zuerst zwei eigene Neuronale Netze gebaut, um die Anwendung der Keras-Bibliothek zu erlernen und den Gebrauch einiger Hilfsfunktionen zu verstehen.
Auf Grundlage der Gruppenergebnisse entschied ich mich bei allen weiteren Netzen für die Verwendung der Aktivierungsfunktion "relu" und des Optimierers "Adam" (vgl. Logbuch von Felix H.).
Ich orientierte mich auch bei Hyperparameterwahl und Struktur des Netzes, d.h. kleiner werdenden Schichten, an den bereits erstellten Modellen der Gruppe. Zudem wurde aus den Gruppenergebnissen
deutlich, dass nicht alle Features für die Vorhersage relevant sind, weshalb ich mich im Folgenden nur mit den Features "MedInc", "AveOccup", "Latitude" und "Longitude" beschäftigte.
Die Auswahl relevanter Features ist ein zentraler Bestandteil von Data-Mining-Prozessen (Runkler, 2012), was die Vorgehensweise bestätigt.
Mit den Netzen erreichte ich  $R^2$-Werte von 0.7188 und 0.7193, was den Vergleich mit der linearen Regression mit $R^2 = 0.7091$ (vgl. \texttt{01\_Lasso\_Ridge.ipynb}) bereits übertraf.

\subsection{Random-Search}
Nun beginnt der Hauptteil meiner Arbeit. Für alle Hyperparameter war von den Gruppenmitgliedern eine ungefähre optimale Größenordnung vorgegeben worden. Diese habe ich verwendet, um
gute Modellkonfigurationen zu generieren. Sorgen machte ich mir, da bisher die Hyperparameter jeweils einzeln variiert wurden, was dazu führen kann, dass bei einer geordneten Suche nicht die optimale Kombination gefunden wird, da der Algorythmus eher Plateaus findet als optimale Kombinationen.
Dies ist darauf zurückzuführen, dass Hyperparameter häufig Wechselwirkungen aufweisen und daher nicht unabhängig voneinander betrachtet werden sollten (Ertel, 2021).
Daher entschied ich mich für die Random-Search-Methode, bei der die Hyperparameter zufällig aus den vorgegebenen Bereichen ausgewählt werden. Zu Beginn wurden diese gewählt, wie in Abbildung \ref{tab:hyperparameters} dargestellt.
Random Search gilt als effektive Methode zur Hyperparameteroptimierung, da sie im Vergleich zu systematischen Suchverfahren häufig schneller gute Kombinationen findet (vgl. Bergstra \& Bengio, 2012).
\begin{table}[ht]
    \centering
    \begin{tabular}{@{}lcc@{}}
        \toprule
        Hyperparameter & Wertebereich
        \\ \midrule
        Anzahl Schichten & 2, 3 \\
        Anzahl Neuronen pro Schicht & 32, 64, 128 \\
        Lernrate & 0.0001 bis 0.01 \\
        Dropout-Rate & 0.1 bis 0.3 \\
        Batchgröße & 16, 32, 64 \\
        Anzahl Epochen & 50, 100 \\
        L2-Regularisierung & 0.00001 bis 0.001 \\
        Aktivierungsfunktion & ReLU \\
        Optimierer & Adam \\
        EarlyStopping & Ja \\
        \bottomrule
    \end{tabular}
    \\
    \caption{Hyperparameter und Wertebereiche für die Random-Search}
    \label{tab:hyperparameters}
\end{table}
\\
Damit die Rechenzeit im Ramen blieb nutze ich Parallelisierung, womit mehrere Modelle gleichzeitig trainiert werden konnten. 
Aus mehreren hundert Modellen, wählte ich die besten 50 aus, um diese zu analysieren. \begin{table}[ht]
    \centering
    \label{tab:random_search_results}
    \resizebox{\columnwidth}{!}{%
    \begin{tabular}{@{}lcccccccc@{}}
        \toprule
        $R^2$-Valuation & $R^2$-Test & units & depth & learning\_rate & l2\_reg & dropout\_rate & batch\_size & epochs \\
        \midrule
        0.757183 & 0.743349 & 128 & 2 & 0.000343 & 0.000052 & 0.156997 & 16 & 81 \\
        0.751369 & 0.736966 & 128 & 2 & 0.001588 & 0.000024 & 0.163947 & 16 & 66 \\
        0.750405 & 0.736433 & 128 & 2 & 0.000395 & 0.000017 & 0.256285 & 64 & 93 \\
        0.747135 & 0.733957 & 64 & 3 & 0.000941 & 0.000284 & 0.116963 & 16 & 77 \\
        0.746669 & 0.733977 & 128 & 3 & 0.001095 & 0.000016 & 0.195057 & 32 & 61 \\
        \bottomrule
    \end{tabular}%
    }
    \\
    \caption{Ergebnisse der Random-Search: Hyperparameterkonfiguration einiger der besten Modelle. }
\end{table}
\\
Aus den Ergebnissen waren keine klaren Aussagen mit bloßem Auge möglich, weshalb 
ich mich dazu entschied einen Random-Forest-Classifier mit den Hyperparametern zu trainieren und damit die Hyperparameter-Importance (Abb. \ref{fig:hyperparameter_importance}) zu erzeugen.
Außerdem liess ich mir die im sklearn-Modul verfügbaren Partial-Dependence-Plots anzeigen, um die Auswirkungen der Hyperparameter auf die Modellleistung zu verstehen.


\begin{figure}[h]
    \centering
    \begin{subfigure}{0.23\textwidth}
        \centering
        \includegraphics[width=\linewidth]{C:/KI_2026/KI_308/ki-308/results/Hyperparameters_Feature_Importance.png}
        \caption{Feature-Importance der Hyperparameter nach der ersten Iteration}
    \end{subfigure}
    \hfill
    \begin{subfigure}{0.23\textwidth}
        \centering
        \includegraphics[width=\linewidth]{C:/KI_2026/KI_308/ki-308/results/Hyperparameters_Feature_Importance_3.png}
        \caption{Feature-Importance der Hyperparameter nach der dritten Iteration}
    \end{subfigure}
    \caption{Vergleich der Hyperparameter-Importance nach der ersten und dritten Iteration der Random-Search.}
    \label{hyperparameter_importance}
\end{figure}

Mit den Partial-Dependence-Plots (Abb. \ref{pdp}) war es dann möglich, die Auswirkungen der Hyperparameter auf die Modellleistung zu verstehen. Durch ein Iteratives Vorgehen,
bei welchem nach jeder Iteration, d.h. nach der Erstellung von einigen hundert Modellen, die Hyperparameter-Importance analysiert und die Wertebereiche der Hyperparameter angepasst wurden, 
konnte das Modell sukzessive verbessert werden. Allerdings zeigte sich nach der zweiten Iteration schon keine große Veränderung mehr, weshalb ich das Vorgehen nach der dritten Iteration beendete. 
\begin{figure}[ht]
    \centering
    \begin{subfigure}{0.45\textwidth}
        \centering
        \includegraphics[width=\linewidth]{C:/KI_2026/KI_308/ki-308/results/PDP_1.png}
        \caption{Partial-Dependence-Plots nach der ersten Iteration}
    \end{subfigure}
    \hfill
    \begin{subfigure}{0.45\textwidth}
        \centering
        \includegraphics[width=\linewidth]{C:/KI_2026/KI_308/ki-308/results/PDP_3.png}
        \caption{Partial-Dependence-Plots nach der dritten Iteration}
    \end{subfigure}
    \caption{Partial-Dependence-Plots nach der ersten und dritten Random-Search-Iteration.}
    \label{pdp}
\end{figure}

Währenddessen probierte ich auch größere Spannen bei den Hyperparametern aus, um zu sehen, ob es vielleicht doch noch bessere Modelle gibt. Zu nennen ist hier zum Beispiel die Tiefe des Netzes oder eine kleinere Lernrate.
Ziel war es, etwas über den Tellerrand hinauszuschauen, um zu prüfen, ob es vielleicht doch noch bessere Modelle gibt. Es zeigte sich, dass die Tiefe des Netzes von 4 Layern gute Ergebnisse lieferte, weshalb ich diesen Wert in die Random-Search mit aufnahm.
Vorerst nutze ich eine maximale Epochenzahl von 100. Da trotz Early-Stopping aber nach der der zweiten Iteration die Netze noch nicht ausreichend trainiert worden waren, erhöhte ich die maximale Epochenzahl.
Schlussendlich passte ich die Hyperparameterbereiche folgendermaßen an (Abb. \ref{tab:hyperparameters_2}):

\begin{table}[ht]
    \centering
    \begin{tabular}{@{}lcc@{}}
        \toprule
        Hyperparameter & Wertebereich
        \\ \midrule
        Anzahl Schichten & 2, 3, 4 \\
        Anzahl Neuronen pro Schicht & 32, 64, 128 \\
        Lernrate & $10^{-3.3}$ bis $10^{-2}$ \\
        Dropout-Rate & 0.115 bis 0.16 \\
        Batchgröße & 16, 32, 64 \\
        Anzahl Epochen & 500 \\
        L2-Regularisierung & 0.00001 bis 0.001 \\
        Aktivierungsfunktion & ReLU \\
        Optimierer & Adam \\
        EarlyStopping & Ja \\
        \bottomrule
    \end{tabular}
    \\
    \caption{Hyperparameter und Wertebereiche für die Random-Search}
    \label{tab:hyperparameters_2}
\end{table}

Das beste gefundene Netz erreichte einen $R^2$-Wert von 0.752 auf dem Testset, was eine deutliche Verbesserung gegenüber der linearen Regression darstellt.


\subsection{Ensemble}
Nachdem ich 10 meiner besten Hyperparameterkonfigurationen aus der Random-Search-Phase ausgewählt hatte, entschied ich mich, ein Ensemble zu erstellen.
Ensemble-Methoden werden eingesetzt, um durch die Kombination mehrerer Modelle die Generalisierungsfähigkeit zu verbessern (Kruse et al., 2013).
Dieses Ensemble verbesserte den Score nicht, sondern erreichte lediglich einen $R^2$-Wert von 0.7499, was trotzdem besser war als die allermeisten Einzelmodelle.
Um eine Aussage über den gesamten California-Housing-Datensatz zu treffen, trainierte ich das Ensemble auf eben diesen und erreichte einen $R^2$-Wert von 0.788, was eine deutliche Verbesserung gegenüber dem verkleinerten Datensatz darstellte.

\subsection{Feature-Engineering}
Aus dem deutlichen Sprung des $R^2$ durch die Hinzunahme der restlichen Features lag nahe, Feature-Engineering
zu betreiben. Tatsächlich konnte ich durch die Erstellung neuer Features (siehe Tab. \ref{tab:features}) das Ensemble in zwei Versuchen nicht verbessern, sondern verschlechterte es sogar. 
\begin{table}[h]
\centering

\begin{tabular}{|c|c|}
\hline
neues Feature & Definition \\
\hline
rooms per household & AveRooms/AveOccupancy \\
bedrooms per room & AveBedrms/AveRooms \\
population per household & Population/AveOccupancy \\
income per household & MedInc/AveOccupancy \\
\hline
\hline
$R^2$-Wert & 0.65 \\
\hline
\end{tabular}
\caption{zusätzliche Features und deren Definitionen im ersten Versuch des Feature-Engineerings}
\vspace{0.5cm}

\begin{tabular}{|c|c|}
    \hline
    neues Feature & Definition \\
\hline
income $\cdot$ rooms & MedInc $\times$ AveRooms \\
income $\cdot$ population & MedInc $\times$ Population \\
distance to center & $\sqrt{(Latitude - 36)^2 + (Longitude + 119)^2}$ \\
squared Occupancy & AveOccup $\cdot$ AveOccup \\
logarithm of population & $\log(Population)$ \\
\hline
\hline
$R^2$-Wert & 0.68 \\
\hline
\end{tabular}

\caption{zusätzliche Features und deren Definitionen im zweiten Versuch des Feature-Engineerings}
\end{table}

I




\begin{thebibliography}{9}

\bibitem{felix_logbuch}
Felix H.
\textit{Logbuch – KI I Projekt, Gruppe 308}.
Unveröffentlichtes Arbeitsdokument, 2026.

\bibitem{florian_logbuch}
Florian.
\textit{Logbuch – KI I Projekt, Gruppe 308}.
Unveröffentlichtes Arbeitsdokument, 2026.

\bibitem{russell2021}
Stuart Russell and Peter Norvig.
\textit{Artificial Intelligence: A Modern Approach}.
4th Edition, Pearson, 2021.

\bibitem{ertel2021}
Wolfgang Ertel.
\textit{Introduction to Artificial Intelligence}.
3rd Edition, Springer, 2021.

\bibitem{runkler2012}
Thomas Runkler.
\textit{Data Mining: Methoden und Algorithmen intelligenter Datenanalyse}.
Vieweg+Teubner, 2012.

\bibitem{kruse2013}
Rudolf Kruse, Christian Borgelt, Frank Klawonn, Christian Moewes, Matthias Steinbrecher.
\textit{Computational Intelligence: Eine methodische Einführung in Künstliche Neuronale Netze, Evolutionäre Algorithmen und Fuzzy-Systeme}.
Springer, 2013.

\bibitem{bergstra2012}
James Bergstra and Yoshua Bengio.
Random Search for Hyper-Parameter Optimization.
\textit{Journal of Machine Learning Research}, 13:281–305, 2012.


\end{thebibliography}




\end{document}